\section{Periodic unique ground state}\label{sec:periodic_unique_ground_state}

\begin{definition}[Span of the periodic MPS vector]\label{def:mpv_submodule}
    \lean{MPSTensor.mpvSubmodule}
    \leanok
    \uses{def:mpv}
    The subspace spanned by the MPS vector
    $V^{(N)}(A)(\sigma) = \tr(A^{\sigma_0}\cdots A^{\sigma_{N-1}})$.
\end{definition}

\begin{lemma}[The periodic MPS vector is the ground-space map at the identity]
    \label{lem:mpv_eq_ground_space_map_one}
    \lean{MPSTensor.mpv_eq_groundSpaceMap_one}
    \leanok
    \uses{def:ground_space_map, def:mpv}
    For every chain length $N$,
    \[
        V^{(N)}(A) = \Gamma_{N}(\Id).
    \]
\end{lemma}

\begin{lemma}[The periodic MPS vector lies in the local ground space]\label{lem:mpv_mem_ground_space}
    \lean{MPSTensor.mpv_mem_groundSpace}
    \leanok
    \uses{def:ground_space_parent, def:mpv}
    $V^{(N)}(A) \in G_N(A)$ for every $N$, with boundary matrix $X = I$.
\end{lemma}

\begin{proof}
    \leanok
    $V^{(N)}(A)(\sigma) = \tr(A^\sigma) = \tr(A^\sigma \cdot I) = \Gamma_N(I)(\sigma)$.
\end{proof}

\begin{definition}[One-dimensional ground space]\label{def:has_unique_ground_state}
    \lean{MPSTensor.HasUniqueGroundState}
    \leanok
    A finite-dimensional ground space $S$ is one-dimensional if $\dim S = 1$.
\end{definition}

\begin{theorem}[One-dimensionality criterion]\label{thm:has_unique_ground_state_iff_proportional}
    \lean{MPSTensor.hasUniqueGroundState_iff_proportional}
    \leanok
    \uses{def:has_unique_ground_state}
    For finite-dimensional $S$, $\dim S = 1$ iff there exists a nonzero
    $\psi_0 \in S$ such that every $\psi \in S$ is of the form $c\,\psi_0$.
\end{theorem}

\begin{proof}
    \leanok
    Apply the standard finite-dimensional criterion that a nonzero space has
    dimension one exactly when all of its vectors are proportional to a single
    nonzero vector.
\end{proof}

\begin{definition}[Periodic parent-Hamiltonian ground space]\label{def:chain_ground_space}
    \lean{MPSTensor.chainGroundSpace}
    \leanok
    \uses{def:ground_space_parent}
    For $N > 0$ and $L \le N$, define the \emph{periodic parent-Hamiltonian
        ground space}
    \[
        \mathcal{G}_{N,L}(A) := \bigl\{\psi \in (\{0,\ldots,d{-}1\}^N \to \C)
        : \forall i \in \{0,\ldots,N{-}1\},\
        \psi|_{[i,i+L-1]} \in G_L(A)\bigr\},
    \]
    where the condition $\psi|_{[i,i+L-1]} \in G_L(A)$ means that for every
    choice of physical indices outside the window, the restriction of $\psi$
    to that window lies in $G_L(A)$.
    The one-index notation $G_L(A)$ denotes the local ground space
    $\mathcal G_L$ of \cite[Section~IV.C]{Cirac2021Matrix}. The two-index
    notation $\mathcal G_{N,L}(A)$ used here denotes the periodic intersection
    of those local constraints over all translated length-$L$ windows on the
    $N$-site ring.
    When $N = 0$ or $L > N$, set $\mathcal{G}_{N,L}(A) := \top$ by convention.
\end{definition}

\begin{lemma}[Cyclic-window witnesses inside the periodic ground space]
    \label{lem:chain_ground_space_window_witnesses}
    \lean{MPSTensor.chainGroundSpace_window_witnesses}
    \leanok
    \uses{def:chain_ground_space, def:ground_space_parent}
    If $0<N$, $L\le N$, and $\psi\in\mathcal G_{N,L}(A)$, then for every
    cyclic window and every outside configuration there is a boundary matrix $Y$
    such that the corresponding restricted $L$-site state is $\Gamma_L(Y)$.
\end{lemma}

\begin{proof}
    \leanok
    This is the defining cyclic-window condition for membership in
    $\mathcal G_{N,L}(A)$, together with the definition of the local ground
    space as the range of $\Gamma_L$.
\end{proof}

\begin{lemma}[The periodic MPS vector satisfies every cyclic window constraint]
    \label{lem:mpv_mem_chain_ground_space}
    \lean{MPSTensor.mpv_mem_chainGroundSpace}
    \leanok
    \uses{def:chain_ground_space, lem:mpv_window_mem_ground_space}
    If $0 < N$ and $L \le N$, then
    \[
        V^{(N)}(A) \in \mathcal G_{N,L}(A).
    \]
\end{lemma}

\begin{theorem}[Periodic MPS vector nonvanishing for block-injective tensors]%
    \label{thm:mpv_ne_zero_blk_injective}
    \lean{MPSTensor.mpv_ne_zero_of_isNBlkInjective}
    \leanok
    \uses{def:mpv, def:l_blk_injective}
    If $A$ is $L_0$-block-injective with $L_0 > 0$, $D \ge 1$, and
    $N \ge L_0 + 1$, then $V^{(N)}(A) \ne 0$ in the $N$-site space.
\end{theorem}

\begin{proof}
    \leanok
    Suppose for contradiction that $V^{(N)}(A) = 0$, i.e.\
    $\tr(A^\sigma) = 0$ for every word $\sigma$ of length~$N$.
    For any words $w$ and $u$ with $|w| = N - L_0$ and $|u| = L_0$,
    we have $\tr(A^w A^u) = 0$.
    Since $A$ is $L_0$-block-injective, $\{A^u : |u| = L_0\}$
    spans $\MN{D}$. By nondegeneracy of the trace pairing,
    $A^w = 0$ for every $|w| = N - L_0$.
    Repeating the argument with words of length $N - 2L_0$ (since
    $|w'| + L_0 = N - L_0$, the product $A^{w'} A^{u_1}$ has length
    $N - L_0$, so $A^{w'} A^{u_1} = 0$ by the previous step) and iterating,
    we eventually reach $\tr(A^u) = 0$ for $|u| = L_0$, which gives
    $I = 0$, a contradiction.
\end{proof}
